% Pre-written reviewer-defense paragraphs for each major brutal-audit finding.
% Not for inclusion in the main paper text — used as raw material for the
% rebuttal letter / discussion section as needed.

\paragraph{Dysarthric-subgroup $\rho = 0.394$ non-significant ($p = 0.33$).}
The $\rho = 0.394$ figure is the dysarthric-subgroup Spearman (the rank
correlation over the 8 dysarthric speakers), \emph{not} a
leave-one-speaker-out cross-validation score. We report it transparently
as a directional indicator only. The 95\% bootstrap CI $[-0.29,\,0.90]$
correctly conveys the uncertainty at $n = 8$; pairwise rank concordance is
$\approx 73\%$. We do not advance it as evidence of cross-speaker
generalisation; a definitive cross-validated test requires $n \geq 20$,
which we propose as the explicit goal of the longitudinal pilot in
Section~6.

\paragraph{MFCC trajectory $\rho$ exceeds ATD globally ($0.849$ vs $0.756$).}
On global Spearman, MFCC--DTW marginally exceeds ATD; the
dysarthric-subgroup $\rho$ for MFCC ($0.282$) is lower than ATD-All's
($0.394$), but both are $n = 8$ non-significant statistics and we read the
difference as directional, not decisive. We do \emph{not} lean on the
controlled-perturbation study to differentiate the two on articulatory
grounds: that study's pre-registered specificity hypothesis was falsified
(pink noise outranks formant shift; ATD does not beat MFCC on formant
shift), and we say so in the Limitations section. The defensible
distinction is that ATD detects frame dropout (partial muting) better than
MFCC, i.e.\ it exploits temporal trajectory structure.

\paragraph{UASpeech null result ($\rho = -0.22$, $p = 0.43$).}
This is reported as a falsified hypothesis, not hidden. The
trajectory-length analysis (TORGO mean $155$ frames vs UASpeech mean $127$)
provides a mechanistic, falsifiable explanation: ATD requires a minimum
temporal regime for DTW signal stabilisation. Track~B replicates the test
on length-stitched UASpeech sentences to test this directly; we
pre-register the prediction that $\rho$ recovers to $\geq 0.4$ on stitched
stimuli. Cross-corpus extension via a shared TORGO PCA basis with
per-corpus CMVN normalisation (Section~5.5) is reported with both
within-corpus and zero-shot transfer numbers.

\paragraph{F04 leverage (removing $1$ speaker shifts $\rho$ from $0.756$ to $0.855$).}
Per Cook's distance, F04 is a high-leverage but not anomalous point --
her clinically documented preserved intelligibility despite a categorical
``severe'' label is itself the substantive finding (Section~6: ATD captures
execution quality independent of categorical labels). We report sensitivity
analyses for every speaker dropped (Table~S1) and the headline result is
robust to F04 inclusion: leave-one-speaker-out across all $15$ speakers
yields $\rho \in [0.71,\,0.86]$ (median $0.76$).

\paragraph{No multiple-comparison correction across $9$ layer tests.}
We apply Benjamini--Hochberg FDR at $q = 0.05$ across all $9$ layers and
$6$ ablation variants ($15$ tests). All $9$ layers remain significant after
correction (smallest adjusted $p = 0.018$); the no-DTW and mean-pool
ablations remain significantly worse than full ATD. Raw and adjusted
$p$-values are reported in Table~2.
